\section{Preparation, whitening, and the smooth limit}
\label{sec:completion}

The finite prediction result encodes a target observable in a packet
sum while retaining nearly flat total spectral density. We prepare
arbitrary observables for that result, whiten the resulting vector to
exactly flat density, and diagonalize the stages to obtain a smooth map
whose one orbit is complete.

All vector comparisons in this section take place in the fixed physical
space $L^2(\T^3)$, even when intermediate coordinates change.  A neighborhood
of a diffeomorphism will mean a neighborhood of the pair consisting of the
map and its inverse in the $C^\infty$ topology.  For a unitary $U$, write
\[
 \calC(U,f)=\overline{\operatorname{span}\{U^nf:n\in\Z\}}.
\]
We use the packet spectral calculus from Section~\ref{sec:passages}:
for a measurable function $\alpha$ of the additive spectral variable $z$, the
coefficient of $\alpha(U_S)F$ at angular index $m$ is
$\alpha(m\phi(u)\bmod1)A_m(u,a)$.  The notation is used only when the resulting
vector has finite $L^2$ norm.

\begin{lemma}[A flat vector]\label{lem:flat-vector}
For every intermediate twist and every good arc $I$, there is a real
centered vector $v$ of spectral density $1$, using two nonzero angular
indices and coefficients supported in a compact subinterval of $I$.
\end{lemma}
\begin{proof}
Choose a positive integer $m$ and $u_0<u_1$ in $I$ such that
$m(\phi(u_1)-\phi(u_0))=1$.  This is possible by first choosing a compact
interval on which $\phi$ increases and then making $m$ large.  Put
\[
 A(u)=\1_{[u_0,u_1]}(u)
       \sqrt{\frac{|m\phi'(u)|}{2}},
 \qquad
 v(\theta,u,a)=A(u)\ee(m\theta)+A(u)\ee(-m\theta).
\]
The pushforward of $|A(u)|^2\dd u\,\dd a$ by each of $\pm m\phi(u)$
modulo one is $\tfrac12\dd z$.  The two angular modes are orthogonal,
so the spectral density of $v$ is $1$.  The vector is real and centered.
Its endpoint discontinuities are harmless: no smoothness is required of
the vector.
\end{proof}

We call a pair $(S,f)$ admissible if $S$ is an intermediate twist in
smooth volume-preserving coordinates, with good arc $I$, and $f$ is real,
has spectral density $1$, and has finitely many nonzero angular modes
with $L^2$ coefficients supported in a compact subset of $I\times\T$.
In particular, $f$ is centered and has norm one.

\begin{proposition}[One stage]\label{prop:stage}
Let $(S,f)$ be admissible, let $h$ be a real smooth function in physical
coordinates with $\int h=0$ and $\|h\|_\infty\le1$, and let
$\eta,\rho,\tau>0$.  In every prescribed neighborhood of $S$ there is an
admissible pair $(S',f')$, with good arc $I'$, and a finite Laurent
polynomial $P$ such that
\begin{equation}\label{eq:stage-conclusion}
 \|f'-f\|_2<\tau,\qquad
 |\T\setminus I'|<\rho,\qquad
 \|h-P(U_{S'})f'\|_2
 <|\T\setminus I|^{1/2}+\eta.
\end{equation}
\end{proposition}

\begin{proof}
We first prepare smooth packets for $f$ and $h$, then put constant tags
on disjoint packets, and finally use Proposition~\ref{prop:prediction}.
Every accuracy appearing below may be made arbitrarily small after the
finite data on which it depends have been fixed.

At a fixed map we use the density estimate
\eqref{eq:density-norm-bound}, which also applies to the nonsmooth
$L^2$ coefficients of an admissible vector. Across a passage, density
preservation instead comes from Lemma~\ref{lem:neutral}.

\paragraph{Converting the angular averages.}
In the current coordinates write
\[
 h=h_\perp+g(u,a),\qquad
 g_0(a)=\int_\T g(u,a)\dd u,\qquad r(u,a)=g(u,a)-g_0(a).
\]
Here $h_\perp$ has zero $\theta$-average, $r$ has zero $u$-average
for each $a$, and $\int g_0=0$.  Fourier truncation and smoothing on the
good arc give a real smooth packet sum $F$ arbitrarily close to $f$.
By \eqref{eq:density-norm-bound}, its density is arbitrarily close to $1$
in $L^1$.

Let $\beta=|\T\setminus I|$.  Fiberwise Parseval gives
\[
 \int_\T|h_\perp(\theta,u,a)|^2\dd\theta
 \le \int_\T|h(\theta,u,a)|^2\dd\theta\le1.
\]
Thus cutting $h_\perp$ off smoothly inside $I$, and then retaining
finitely many of its nonzero angular modes, loses at most
$\sqrt\beta+\eps$ in norm, for arbitrary $\eps>0$.  For example, the
cutoff may be chosen with
$\int_\T(1-\chi)^2\dd u<\beta+\eps^2$.
Denote the resulting packet sum by $H_\perp$.
The simultaneous-list version of Lemma~\ref{lem:neutral}, stated in
Remark~\ref{rem:finite-family-passages}, transports the $F$ and
$H_\perp$ lists through a common passage even when their supports
overlap.  We require density preservation only for the $F$ list; no
orthogonality between the two lists is asserted.

The next passage turns variation in the old $u$ coordinate into
nonzero angular modes; the passage after it does the same for the old
$a$ coordinate. There are two remaining terms, and neither will be
discarded. Before
the first passage integer $N_1$ is chosen, fix a finite nonzero
$u$-Fourier truncation $r_K$ of $r$.  Choose the new good arc in the
$a$ variable and its smooth cutoff $\chi_1(a)$ so that replacing $r_K$
by $\chi_1r_K$ costs arbitrarily little.  These are fixed functions of
the old slow variables.  The cutoff can simultaneously lose arbitrarily
little mass from the finitely many existing packets.  Choose the broad
neutral jitter and then $N_1$, as in Lemma~\ref{lem:neutral}.  In the new
coordinates
\[
 (u,a,b),\qquad b=\theta-H_1(u),
\]
the function $\chi_1(a)r_K(u,a)$ is already a finite sum of nonzero
angular modes supported in the new good arc.  The former packets are
transported with arbitrarily small norm error, and the density of the
$F$ packets changes by arbitrarily little in $L^1$.  Choose $N_1$ also so that every retained transported old index lies
outside $[-K,K]$: its quotient by $N_1$ stays in a fixed compact set
away from zero. This keeps the added modes of $r_K$ disjoint in angular
index from the transported test packets.  The function $g_0(a)$ has been left entirely untouched.

Now fix a finite nonzero $a$-Fourier truncation $(g_0)_L$ of $g_0$.
It has no constant term because $\int g_0=0$.  Apply a second neutral
passage to all existing packets.  Its new coordinates are
\[
 (a,b,d),\qquad d=u-H_2(a).
\]
Choose its good arc and cutoff $\chi_2(b)$ before its final integer
$N_2$.  The term $\chi_2(b)(g_0)_L(a)$ is now a smooth packet sum in
nonzero angular modes.  Its cutoff loses arbitrarily little norm and
cannot introduce an angular constant.  Choose $N_2$ also so that the transported indices lie outside
$[-L,L]$, keeping the newly added modes disjoint from them.  The existing packets are
transported with arbitrarily small norm and density errors.

Relabel the coordinates of the resulting twist as $(\theta,u,a)$.
For every $\eps>0$ the above choices therefore give real smooth packet
sums $F,H$ such that
\begin{equation}\label{eq:prepared-pair}
 \|F-f\|_2<\eps,\qquad
 \|H-h\|_2<\sqrt\beta+\eps,\qquad
 \|p_F-1\|_1<\eps.
\end{equation}
The functions that needed new cutoffs were fixed before the relevant
passage integer was chosen.  In particular, this argument has not
assumed uniform control of coefficients arising only after that choice.

\paragraph{Constant ratios on rectangles.}
Enlarge the finite angular index list so that it contains the modes of
both $F$ and $H$, and write their coefficients as $F_m,H_m$.  For almost
every sufficiently small real $t$,
\begin{equation}\label{eq:ratio-null-set}
 \int_{\{F_m+tH_m=0\}}|H_m|^2\dd u\,\dd a=0
 \quad\hbox{for every }m.
\end{equation}
To see this, at a point where $H_m\ne0$ at most one real $t$ solves
$F_m+tH_m=0$; Fubini applied to $|H_m|^2\dd u\,\dd a\,\dd t$ proves
the assertion.  There are only finitely many indices.  Replace $F$ by
$F+tH$ with such a sufficiently small real $t$.  Its norm and density
changes are arbitrarily small by \eqref{eq:density-norm-bound}, and
reality is preserved.

Define $R_m=H_m/F_m$ where $F_m\ne0$, and zero otherwise.  In view of
\eqref{eq:ratio-null-set}, $H_m=R_mF_m$ almost everywhere.  With
$\dd\nu_m=|F_m|^2\dd u\,\dd a$, one has
$\int|R_m|^2\dd\nu_m=\|H_m\|_2^2$.  Truncating $R_m$ to a bounded
function therefore changes $R_mF_m$ arbitrarily little in $L^2$.
Finite rectangular step functions are dense in $L^2(\nu_m)$: continuous
functions are dense for this finite measure, and uniform continuity
approximates a continuous function by a fine rectangular grid.  Grid
boundaries have $\nu_m$-measure zero, since $\nu_m$ is absolutely
continuous.

Approximate each bounded ratio by such a step function, and replace the
indicators of the finitely many grid rectangles by smooth interior
cutoffs.  Trim enough to obtain positive gaps between distinct retained
rectangles, but lose arbitrarily little in both the $F$ norm and the
tagged norm.  This is possible after the finitely many step values are
known.  Where a ratio has been truncated, a zero step value is allowed;
we do not need to remove the corresponding $F$ mass.  Performing the
construction at positive indices and reflecting it at negative indices
gives packets $F_r$ and constants $b_r$ with conjugate symmetry, such
that
\[
 F^{\mathrm{tag}}=\sum_r F_r\quad\hbox{approximates }f,
 \qquad
 H^{\mathrm{tag}}=\sum_r b_rF_r\quad\hbox{approximates }h
\]
with the errors in \eqref{eq:prepared-pair}, up to arbitrarily small
additions.  Different rectangles in $(m,u,a)$ are disjoint with gaps;
each is its own color.  The density of $F^{\mathrm{tag}}$ remains
arbitrarily close to $1$ in $L^1$.  The number
\[
 M=\max\bigl(1,\max_r|b_r|\bigr)
\]
is finite.  All subsequent approximation errors for tagged vectors are
chosen after $M$ is known.

\paragraph{Separation by angular index alone.}
Two further neutral passages supply exactly the hypotheses of
Proposition~\ref{prop:prediction}.  Here are the support details.
Write the single old angular index of rectangle $r$ as $m_r$, and its
compact coefficient support as $K_r$.  Choose $c$ so that different
$m_r$ have separated ranges, and form
\[
 E_r=\{(m_r(c+\phi(u)),a):(u,a)\in K_r\}.
\]
The compact sets $E_r$ are pairwise disjoint: for a fixed $m_r$ the map
$(u,a)\mapsto(m_r(c+\phi(u)),a)$ is injective, and different indices
were separated by $c$.  Choose smooth cutoffs $\xi_r(s,a)$ supported
in disjoint neighborhoods of these sets and equal to one near them.
On the first passage retain the exact coefficients
\begin{equation}\label{eq:first-separation}
 C_{j,r}(a)=\chi(a)\xi_r(j/N,a)
 \int_\T A_r(u,a)\ee(m_rH(u)-ju)\dd u.
\end{equation}
Here $\chi$ is the new good-arc cutoff; its loss was made small before
choosing the broad neutral jitter and $N$.  Lemma~\ref{lem:passage}
shows that the support restrictions delete arbitrarily little norm as
$N$ increases.  There are finitely many retained $j$, all nonzero.  In
coordinates $(u,a,b)$ the resulting packets are
\[
 C_{j,r}(a)\ee(ju)\ee(m_rb).
\]
Their moduli are independent of the new spare coordinate $b$.
Different old indices cannot share $j$; at a shared $j$, distinct colors
have disjoint compact $a$ supports with positive gaps.

Fix this finite output before making the second passage.  Choose its
$c_2$ to separate the ranges belonging to distinct $j$.  For each fixed
$j$, strict monotonicity of the current speed $\sigma$ on its good arc
makes the compact sets
\[
 Q_{j,r}=\{j(c_2+\sigma(a)):a\in\supp C_{j,r}\}
\]
disjoint as $r$ varies.  Thus all nonempty $Q_{j,r}$ have disjoint open
neighborhoods.  Restrict the second-passage coefficient at new index
$k$ by a smooth cutoff of $k/N_2$ supported in its designated
neighborhood and equal to one near $Q_{j,r}$.  Taking $N_2$ last makes
the deleted norm arbitrarily small, while the retained finite
$k$-index sets are \emph{exactly} disjoint for distinct pairs $(j,r)$.
An individual output term has the form
\[
 B_{k,j,r}(b)\ee(ka)\ee(jd),\qquad d=u-H_2(a),
\]
so its modulus is independent of the final spare coordinate $d$.
No interference can spoil this property, because different pairs have
disjoint final index sets.

Every integral and cutoff used here commutes with the constant tag
$b_r$.  Consequently, if the sum of the individual color transport
errors is less than $\varepsilon/(1+M)$, both the untagged and the
tagged physical errors are less than $\varepsilon$.  Neutral density
preservation and \eqref{eq:density-norm-bound} keep the total density
arbitrarily close to $1$.  At this point the packet colors are disjoint
in angular indices, their moduli are independent of the spare
coordinate, and their mass can be required to be at most $3/2$.

\paragraph{Prediction and whitening.}
Apply Proposition~\ref{prop:prediction} to the tags $b_r/M$, with
prediction tolerance divided by $M$.  Its physical transport errors
are likewise chosen with the factor $M$ included.  Write $G$ for the
final real packet sum, $p$ for its density, and $D(U)G$ for the
resulting approximation to the tagged vector, with the factor $M$
restored.  Here $U$ is the Koopman operator of the final intermediate
map.  We can ensure
\begin{equation}\label{eq:prediction-ready}
 \|G-f\|_2\longrightarrow0,
 \qquad
 \|h-D(U)G\|_2<\sqrt\beta+\eta/2,
 \qquad
 \int_\T\sqrt p\,\dd z\longrightarrow1
\end{equation}
as the preparation and prediction tolerances tend to zero.  For the
last assertion, if the input density is $p_0$, then
\[
 \left|\int\sqrt{p_0}-1\right|
 \le\|\sqrt{p_0}-1\|_2
 \le\|p_0-1\|_1^{1/2},
\]
and Proposition~\ref{prop:prediction} gives an arbitrarily small loss
in $\int\sqrt p$.  Its upper limit is at most one by
$\int\sqrt p\le\|G\|_2$ and $\|G-f\|_2\to0$.

Reality implies $p(-z)=p(z)$ almost everywhere.  Set $E=\{p=0\}$,
choose the flat real vector $v$ of Lemma~\ref{lem:flat-vector} for this
last map, and define
\begin{equation}\label{eq:whitening}
 f'=\left(\frac{\1_{E^c}}{\sqrt p}\right)(U)G+\1_E(U)v.
\end{equation}
The first term is well defined in $L^2$, with squared norm $|E^c|$.
The terms have disjoint spectral supports and densities $\1_{E^c}$
and $\1_E$, respectively.  Indeed, $\1_E(U)G=0$, even if $v$ uses some
of the same angular modes.  The total density of $f'$ is exactly one.
Both multipliers preserve reality, nonzero angular indices, and compact
support in the good arc.  Hence the resulting pair is admissible.
Directly, including the contribution from $E$,
\begin{equation}\label{eq:whitening-distance}
 \|f'-G\|_2^2
 =\int_\T(\sqrt p-1)^2\dd z
 =\|G\|_2^2+1-2\int_\T\sqrt p\,\dd z.
\end{equation}
This tends to zero by \eqref{eq:prediction-ready}, and therefore
$\|f'-f\|_2$ can be made less than the independently specified $\tau$.

Whitening also preserves the prediction.  In the measurable spectral
calculus,
\[
 G=(\sqrt p)(U)f',\qquad
 D(U)G=(D\sqrt p)(U)f'.
\]
Since $f'$ has density one and $D(U)G\in L^2$, the function
$D\sqrt p$ belongs to $L^2(\T,\dd z)$.  Trigonometric polynomials are
dense in this space, so $D(U)G$ belongs to $\calC(U,f')$ and is
approximated by a finite Laurent polynomial in $U$ applied to $f'$.
Choose that approximation with error less than $\eta/2$.  This proves
the last inequality in \eqref{eq:stage-conclusion}; it asserts an
approximation to $h$, not exact membership of $h$ at this finite stage.

Finally, each neutral passage and Proposition~\ref{prop:prediction}
allows an arbitrarily small physical perturbation of both map and
inverse.  There are finitely many such operations here, so they can
all be kept in the prescribed neighborhood.  The last new good arc
may have complement less than $\rho$.  All choices of small physical
errors are made after the finite tags, cutoffs, and coordinate changes
that affect their bounds have been fixed.  Thus neither large tags nor
large derivatives of an earlier coordinate change obstruct the
arbitrary closeness claimed in the proposition.
\end{proof}

\begin{proof}[Proof of Theorem~\ref{thm:main}]
Choose a smooth periodic speed $\phi_0$ which agrees with an increasing
affine function on a nonempty closed interval inside $(0,1)$, and let
$I_0$ be an open subinterval where $\phi_0'>0$.  Start with the twist
$S_0$ and its flat vector $f_0$ from Lemma~\ref{lem:flat-vector}.
Choose a countable family of real smooth centered functions, bounded
by one, whose real linear span is dense in real $L^2_0(\T^3)$; the
physical torus sines and cosines will do.  List them as $(h_n)_{n\ge1}$
so that each occurs infinitely often.

We record the continuity needed to retain approximations.  If
volume-preserving maps $T_i$ and their inverses converge uniformly to
$T$ and $T^{-1}$, then for every fixed $k\in\Z$,
\begin{equation}\label{eq:strong-koopman}
 U_{T_i}^kg\longrightarrow U_T^kg\quad\hbox{in }L^2
 \quad(g\in L^2).
\end{equation}
Uniform convergence of the corresponding fixed iterates follows by
composition on the compact torus.  For continuous $g$ this proves
uniform convergence of the compositions; for general $g$, approximation
by continuous functions and unitarity prove \eqref{eq:strong-koopman}.
Consequently, for every fixed Laurent polynomial $P$ the map
\[
 (T,g)\longmapsto P(U_T)g
\]
is continuous into physical $L^2$ when maps and inverses have the
$C^\infty$ topology and vectors have the $L^2$ topology.  Explicitly,
if $P(w)=\sum_k c_kw^k$, its difference is bounded by
\[
 \sum_k|c_k|\|g_i-g\|_2+
 \sum_k|c_k|\|U_{T_i}^kg-U_T^kg\|_2.
\]

Inductively apply Proposition~\ref{prop:stage} with $h=h_n$ and
$\eta=2^{-n}$, obtaining $(S_n,f_n)$, a good arc $I_n$, and a finite
Laurent polynomial $P_n$, such that
\begin{equation}\label{eq:stage-polynomial}
 |\T\setminus I_n|<4^{-n},\qquad
 \|h_n-P_n(U_{S_n})f_n\|_2
 <|\T\setminus I_{n-1}|^{1/2}+2^{-n}.
\end{equation}
Use the arbitrary closeness in that proposition also to impose
\[
 \|f_n-f_{n-1}\|_2<2^{-n},
 \qquad
 d_k(S_n,S_{n-1})+d_k(S_n^{-1},S_{n-1}^{-1})<2^{-n}
 \quad(0\le k\le n),
\]
where $d_k$ are ambient $C^k$ sup-norm distances induced by a fixed
smooth embedding of the torus. At the same time impose
the finitely many open conditions
\begin{equation}\label{eq:retain-polynomial}
 \|P_j(U_{S_n})f_n-P_j(U_{S_{n-1}})f_{n-1}\|_2<2^{-n}
 \qquad(1\le j<n).
\end{equation}
Their compatibility follows from the continuity just proved.  No bound
uniform in $j$ on the degrees or coefficients of $P_j$ is required.

For each fixed $k$, the map and inverse sequences have summable tails
in $C^k$.  They therefore converge smoothly to maps $T$ and $R$.
Uniform convergence, compactness, and
$S_nS_n^{-1}=S_n^{-1}S_n=\operatorname{id}$ give
$TR=RT=\operatorname{id}$, so $T$ is a smooth diffeomorphism with
smooth inverse.  For every continuous $g$,
\[
 \int g\circ T=\lim_n\int g\circ S_n=\int g,
\]
so normalized volume is invariant.  The vectors $f_n$ converge in
$L^2$ to a real centered vector $f$.

Every completed stage has density exactly one.  Combining
\eqref{eq:strong-koopman} with $f_n\to f$, for all $r,s\in\Z$ we obtain
\begin{equation}\label{eq:limit-orthonormal}
 \langle U_T^rf,U_T^sf\rangle
 =\lim_n\langle U_{S_n}^rf_n,U_{S_n}^sf_n\rangle
 =\delta_{rs}.
\end{equation}
The retained finite polynomials give, for every $j$,
\[
 \|h_j-P_j(U_T)f\|_2
 \le |\T\setminus I_{j-1}|^{1/2}+2^{-j}
       +\sum_{n>j}2^{-n}
 =|\T\setminus I_{j-1}|^{1/2}+2^{1-j}.
\]
Along the infinitely many occurrences of any fixed test the right side
tends to zero.  Hence every test belongs to $\calC(U_T,f)$.  Their
complex linear span is dense in complex $L^2_0(\T^3)$, whereas every
orbit vector is centered.  Thus
\[
 \calC(U_T,f)=L^2_0(\T^3).
\]
Together with \eqref{eq:limit-orthonormal}, this says that the bilateral
orbit of $f$ is an orthonormal basis of the entire centered space.
With $w=\ee(z)$, the assignment $U_T^nf\mapsto w^n$ extends to a
unitary map onto $L^2(S^1,m)$ intertwining $U_T$ with multiplication
by $w$.

Multiplication by $w$ has no nonzero invariant vector, since
$(w-1)g(w)=0$ forces $g=0$ almost everywhere.  Consequently every
$T$-invariant $L^2$ function is constant, which proves ergodicity.
The manifold $\T^3$ is compact, connected, smooth, and without
boundary, and its normalized flat volume has strictly positive smooth
density.  These are all the properties required in
Theorem~\ref{thm:main}.
\end{proof}
